\section{Validation Strategy Overview}

RunLine is a web application that generates running loops. A runner chooses a
start, a target distance, a surface preference (road, mixed or trail) and an
elevation preference (flat, balanced or hilly), and RunLine returns up to three
loops ranked by how well they fit the request. Each route can be downloaded as
GPX for a watch or navigation app. The routing engine searches an
OpenStreetMap walking network \citep{boeing2017osmnx,hagberg2008networkx} with
a cost function that penalizes mapped traffic signals, major-road crossings,
repeated segments, the wrong surface and the wrong amount of climbing.

For this report, ``working'' means what our Project Proposal \& Plan
\citep{runguys2026proposal} said it would mean. The proposal defined success as a
runner being able to enter a start and a distance, receive several routes that
stay close to that distance, compare measurable tradeoffs between them, and
export the one they choose. It turned that goal into eleven user stories with
numeric acceptance criteria: seven must-haves (US1 to US7), three
should-haves (US8 to US10) and one nice-to-have (US11). We use those criteria
unchanged as the definition of ``working'' in this report, because they were
fixed before most of the system was built and so cannot have been tuned to
the results. Table~\ref{tab:criteria} restates them. The proposal also set two
technical questions: whether interruption-aware scoring reduces mapped
signals and crossings without losing distance accuracy (US3), and whether the
surface and elevation preferences change the selected routes in the intended
direction (US5, US6).

Two further expectations came from the proposal's discussion of limits rather
than from a user story, and we hold the system to them as well. RunLine must
present route characteristics, never safety guarantees, since OpenStreetMap
can omit crossings, signals and restrictions. It must also avoid retaining
runners' locations, because a start is often a home.

\begin{table}[ht]
\centering\small
\caption{Acceptance criteria from the Project Proposal \& Plan, used as the definition of success.}
\label{tab:criteria}
\begin{tabularx}{\linewidth}{@{}l l Y@{}}
\toprule
ID & Priority & Acceptance criterion (abridged) \\
\midrule
US1 & Must & $\geq$80\% of successful benchmark requests return a top-ranked route within 5\% of the requested distance; all return a route within 10\%. \\
US2 & Must & For $\geq$80\% of requests where the network permits it, at least two candidates overlap by no more than 80\%. \\
US3 & Must & Against a distance-only baseline, median signals plus major-road crossings fall by $\geq$20\% with median distance error $\leq$5\%. \\
US4 & Must & Each candidate shows distance error and the major score components, matching the calculated metrics within display rounding. \\
US5 & Must & Where more than one surface exists, the route moves toward the requested surface relative to the neutral setting in $\geq$75\% of cases. \\
US6 & Must & The hilly route has more elevation gain than the flat route in $\geq$80\% of paired requests. \\
US7 & Must & Ten sampled routes produce GPX that parses without error; one route per distance opens in a GPX-compatible application. \\
US8 & Should & All three team members complete a browser checklist (search, map click, generate, switch, export) without blocking errors. \\
US9 & Should & With a drive radius set, every suggested start lies within the radius and the comparison shows route score and driving cost; ten cases checked. \\
US10 & Should & Warm 3 and 5 mile requests: median below 10\,s, 95th percentile below 20\,s. \\
US11 & Nice & At least one additional area, built with the same pipeline, passes the distance, validity, browser and GPX smoke tests. \\
\bottomrule
\end{tabularx}
\end{table}

To decide whether RunLine meets those criteria we used five validation
methods, each chosen for what it can establish about a route planner and what
the others cannot. Appendix~\ref{app:trace} traces every criterion through the
method that tested it, the artifact that ran, the evidence it produced and the
decision we took.

\begin{enumerate}
\item \textbf{Automated tests.} A pytest suite and two Node test files check
  the deterministic pieces: distance and overlap calculations, signal and
  crossing counts, repeated-segment detection, scoring, GPX generation, pool
  storage and serving, coverage errors, and the browser's request and
  recording logic. These establish that each component computes what it
  claims to on known inputs. They cannot show that the routes are good.
\item \textbf{A fixed route benchmark measured against the acceptance
  criteria.} Eight public Charlottesville starts at 3, 5, 8 and 12 miles
  (32 requests) are planned under each preference setting, and a script
  computes every numeric criterion (US1, US2, US3, US5, US6, US7 and the local
  part of US10). This is the method the proposal committed to. Because the
  inputs are fixed, it also works as a regression test: our repository rules
  require before-and-after benchmark and shape-audit numbers for any change
  to the routing, scoring or model code before it can be merged.
\item \textbf{An independent audit of the crossing metric.} US3 depends on
  the planner counting major-road crossings correctly, and a metric that only
  agrees with itself proves nothing. We built a separate geometric labeller
  that cuts each route's corridor along major-road centrelines and counts a
  crossing whenever the route moves from one side to the other. We then
  checked it against hand-labelled cases and measured the planner's count
  against its labels.
\item \textbf{Validation of the deployed system.} Scripted requests against
  the public URL check what a runner actually receives: response times
  (US10), coverage in the additional metro areas (US11), closed loops and
  matching GPX, error handling, nearby-start suggestions (US9), and whether
  the numbers the interface displays match the API's metrics (US4).
\item \textbf{Manual review by the team.} The browser checklist (US8), the
  ten nearby-start cases (US9), opening GPX files in an outside application
  (US7) and a three-person route review cover what scripts cannot judge:
  whether the interface is usable and whether a route looks like something a
  runner would actually run.
\end{enumerate}

We deliberately did not run a study with outside participants, as the
proposal stated. Recruiting runners would have required consent procedures
and handling of participants' locations. Our acceptance criteria are also
mostly properties of routes, which a fixed benchmark measures more
repeatably than a small participant group would. The cost is that no
evidence in this report shows that runners other than ourselves find the
routes or the score breakdown useful. We return to this gap in
Sections~\ref{sec:analysis} and~\ref{sec:readiness}.

The evidence covers a specific version of the system. The benchmark measured
\texttt{main} at commit \texttt{f482881} \citep{runline2026repo}, and the
deployed checks measured the production site at
\url{https://runline-nine.vercel.app}. During the first deployed run on
16 November, the production backend reported commit \texttt{ebcd150} from
11 November, three code pull requests behind \texttt{main}. None of them
changed scoring, but one (\#26) changed how loops are generated in rare,
degenerate cases. The backend was redeployed on 18 November at
\texttt{62e37f5}, which includes every change the benchmark measured plus the
fixes described in Section~\ref{sec:analysis}, and the deployed checks were
rerun against it. Charlottesville
is the primary area. Its graph was rebuilt in November to keep signalized
crosswalks, and it is served with a complete set of precomputed routes
(11,417 pools, 605,206 loops). Washington DC, New York, Los Angeles, Atlanta,
Boston and Phoenix are served from a team member's virtual machine and are
covered only as US11 evidence. The stakeholders are the three team members
acting as runners and the course staff. The features covered are loop
generation, preferences, ranking, the comparison interface, GPX export,
nearby starts and start-to-finish routes.

We could not validate several things. In-app GPS run recording was added late
and has only been tested on synthetic location samples, never outdoors on a
phone. Accessibility was checked only by an automated audit; we did not test
keyboard-only use or a screen reader. We did not load-test the deployment with concurrent users. Washington,
New York and Los Angeles were built from an older map graph that does not
keep signalized crosswalks, so their signal counts are less complete;
Atlanta, Boston and Phoenix were built with the signal-preserving builder. We also cannot independently verify OpenStreetMap's tags in the
field, so every interruption count in this report means \emph{mapped}
signals and crossings \citep{osmwiki2026signals}.
